% titel: Blatt 6, A2: Lewis-Strukturformeln zeichnen

\subsection*{Lösung (Beschreibung der wesentlichen Strukturmerkmale)}
Sämtliche Formeln sind nach dem Schema Zeichnen --- hier die tragenden
Strukturmerkmale je Molekül:

\begin{itemize}
  \item \textbf{$\mathrm{COS}$:} linear, Zentralatom C; $\mathrm{O=C=S}$
        (zwei Doppelbindungen, keine freien Paare am C).
  \item \textbf{$\mathrm{SeF_4}$:} Se (Periode~4) zentral, 4 Einfachbindungen zu
        F, dazu ein freies Paar am Se (insgesamt 5 Elektronenpaare,
        see-saw-förmig).
  \item \textbf{$\mathrm{SH_3^+}$:} S zentral mit 3 H (Einfachbindungen) und
        einem freien Paar; Kation (eine positive Formalladung am S).
  \item \textbf{$\mathrm{SOCl_2}$:} S zentral, eine $\mathrm{S=O}$-Doppelbindung,
        zwei $\mathrm{S-Cl}$-Einfachbindungen, ein freies Paar am S.
  \item \textbf{$\mathrm{NO_2F}$:} N zentral, eine $\mathrm{N=O}$-Doppelbindung,
        eine $\mathrm{N-O}$-Einfachbindung und eine $\mathrm{N-F}$-Einfachbindung,
        kein freies Paar am N (trigonal-planar).
  \item \textbf{$\mathrm{PCl_6^-}$:} P zentral, 6 $\mathrm{P-Cl}$-Einfachbindungen,
        oktaedrisch, Gesamtladung $-1$ (Oktetterweiterung, P in Periode~3).
  \item \textbf{$\mathrm{BF_4^-}$:} B zentral, 4 $\mathrm{B-F}$-Einfachbindungen,
        tetraedrisch, Ladung $-1$.
  \item \textbf{$\mathrm{AsCl_3}$:} As zentral, 3 $\mathrm{As-Cl}$-Einfachbindungen
        und ein freies Paar (trigonal-pyramidal).
  \item \textbf{$\mathrm{ICl_3}$:} I zentral, 3 $\mathrm{I-Cl}$-Einfachbindungen
        und zwei freie Paare (T-förmig).
  \item \textbf{$\mathrm{H_2CO}$}: C zentral, eine $\mathrm{C=O}$-Doppelbindung,
        zwei $\mathrm{C-H}$-Einfachbindungen, trigonal-planar.
  \item \textbf{$\mathrm{CO}$:} $\mathrm{C\equiv O}$ (Dreifachbindung) und je ein
        freies Paar an C und O.
\end{itemize}

\begin{ergebnis}
Die gezeichneten Strukturen erfüllen alle Oktetts (Periode-2-Zentren) bzw.\
erlauben bei Periode-3-Zentren (P, S, Se, I) die nötige Oktetterweiterung.
\end{ergebnis}

\begin{kontrolle}
Elektronenbilanz stichprobenartig: $\mathrm{CO}$ hat $4+6=10\,\mathrm{e^-}$; eine
$\mathrm{C\equiv O}$-Dreifachbindung (6) plus zwei freie Paare (4) ergibt $10$,
oktett-konform. $\mathrm{PCl_6^-}$: $5+6\cdot 7+1=48\,\mathrm{e^-}$ = 6 Bindungen
(12) + 18 $\mathrm{e^-}$ als freie Paare an den 6 Cl.
\end{kontrolle}

\begin{merksatz}
Zuerst Valenzelektronen zählen und das Zentrum wählen; Einfachbindungen, dann
freie Paare, erst zuletzt Mehrfachbindungen --- und für Periode-3-Zentren ist
Oktetterweiterung erlaubt.
\end{merksatz}
